% % =============================================================
% %  SECTION 3 — METHODOLOGY  (revised)
% %  Added: algorithm pseudocode placeholders, parameter table,
% %         per-module I/O specification, reproducibility subsection.
% %  Removed: generic design rationale prose, repeated figure text.
% % =============================================================
% \section{Methodology}
% \label{sec:methodology}

% \subsection{Problem Formulation}

% Let a multimodal business input be a pair $\mathcal{M} = (V, A)$ where
% $V$ denotes a visual signal (a scanned or photographed document image) and
% $A$ denotes an acoustic signal (a WAV-format spoken purchase order). Either
% component may be absent; the pipeline processes whichever modalities are
% present.

% The overall system implements two composable functions:
% \begin{equation}
%     f_{\mathrm{IE}}: \mathcal{M} \to \mathcal{E}, \qquad
%     f_{\mathrm{ER}}: \mathcal{E} \to \mathcal{D},
% \end{equation}
% where $f_{\mathrm{IE}}$ is the information extraction function,
% $\mathcal{E} = \{e_1, \ldots, e_n\}$ is a set of structured entity tuples,
% and $f_{\mathrm{ER}}$ maps each extracted entity surface form to a canonical
% entry in the enterprise master database $\mathcal{D}$. The entity schema
% captures four attribute types: \textsc{Merchant}, \textsc{Product},
% \textsc{Quantity}, and \textsc{Unit}.

% The modular decomposition of $f_{\mathrm{IE}}$ and $f_{\mathrm{ER}}$ allows
% each component to be tuned or replaced independently, which is important in
% enterprise settings where the product catalog and merchant registry evolve
% continuously.

% \subsection{System Architecture}
% \label{sec:architecture}

% Figure~\ref{fig:architecture} shows the eight-layer pipeline. The system is
% stateless between documents and processes each input as a self-contained batch
% unit, enabling horizontal scaling.

% \begin{figure}[htbp]
% \centering
% \begin{tikzpicture}[
%     box/.style={draw, rectangle, rounded corners, align=center,
%                 minimum width=3.4cm, minimum height=0.9cm, font=\small},
%     smallbox/.style={draw, rectangle, rounded corners, align=center,
%                      minimum width=3.0cm, minimum height=0.75cm, font=\footnotesize},
%     erbox/.style={draw, rectangle, rounded corners, align=center,
%                   minimum width=6.5cm, minimum height=1.5cm, font=\footnotesize},
%     arrow/.style={-Latex, thick},
%     dashedarrow/.style={-Latex, thick, dashed}
% ]
% \node[box] (upload)  at (0, 0)    {Layer 1: Batch Ingestion};
% \node[box] (router)  at (0, -1.5) {Layer 2: \textbf{Modality Router}};
% \node[box] (imgpipe)   at (-4.0, -3.2) {Layer 3: Visual Pipeline};
% \node[box] (audiopipe) at ( 4.0, -3.2) {Layer 3: Acoustic Pipeline};
% \node[smallbox] (ocr) at (-4.0, -4.8) {Layer 4: OCR Transcription};
% \node[smallbox] (asr) at ( 4.0, -4.8) {Layer 4: ASR Transcription};
% \node[box, minimum width=7.5cm] (llm) at (0, -6.5)
%     {Layer 5: LLM-based Entity Extraction ($f_{\mathrm{IE}}$)};
% \node[box] (json) at (0, -8.0)
%     {Layer 6: Unified Structured Entities $\mathcal{E}$};
% \node[erbox, minimum width=8cm] (er) at (0, -10.2)
%     {Layer 7: \textbf{Entity Resolution ($f_{\mathrm{ER}}$)}\\[4pt]
%      Merchant: Fuzzy String Matching\\[2pt]
%      Product: Hybrid Vector Retrieval\\[2pt]
%      (Optional) Taxonomy-Guided Filter};
% \node[smallbox, text width=2.5cm, right=1.8cm of er] (hitl)
%     {Top-$K$\\Human Review};
% \node[box] (output) at (0, -12.5)
%     {Layer 8: Canonical Database Integration $\mathcal{D}$};

% \draw[arrow] (upload)  -- (router);
% \draw[arrow] (router) -| node[pos=0.25, above, font=\scriptsize] {Visual} (imgpipe);
% \draw[arrow] (router) -| node[pos=0.25, above, font=\scriptsize] {Acoustic} (audiopipe);
% \draw[arrow] (imgpipe) -- (ocr);
% \draw[arrow] (audiopipe) -- (asr);
% \draw[arrow] (ocr) |- (llm);
% \draw[arrow] (asr) |- (llm);
% \draw[arrow] (llm) -- (json);
% \draw[arrow] (json) -- (er);
% \draw[arrow] (er)  -- (output);
% \draw[dashedarrow] (er.east) -- node[above, font=\scriptsize] {Low conf.} (hitl.west);
% \draw[dashedarrow] (hitl.south) |- (output.east);

% \node[font=\scriptsize, gray, rotate=90, anchor=south] at (-6.2, -1.5) {Signal Triage};
% \node[font=\scriptsize, gray, rotate=90, anchor=south] at (-6.2, -5.6) {Two-Stage Extraction};
% \node[font=\scriptsize, gray, rotate=90, anchor=south] at (-6.2, -10.2) {Entity Resolution};
% \end{tikzpicture}
% \caption{Architecture of the proposed multimodal information extraction and
% entity resolution pipeline. Layers 1--4 perform signal triage and modality-specific
% transcription; Layers 5--6 extract structured entities via LLM; Layer 7 resolves
% entities against the enterprise catalog; Layer 8 writes canonical records. The dashed
% path routes low-confidence candidates to human review before integration.}
% \label{fig:architecture}
% \end{figure}

% \subsection{Stage 1: Modality Transcription}
% \label{sec:transcription}

% \textbf{Input}: A raw binary file of known MIME type ($V$ or $A$, or both).\\
% \textbf{Output}: A UTF-8 text string $t$ containing the document content,
% with layout structure preserved where available.

% \paragraph{Visual channel.}
% Document images are submitted to a cloud OCR service (Google Document AI)
% configured with the \texttt{DOCUMENT\_TEXT\_DETECTION} processor, which
% returns character-level bounding boxes and a reading-order text assembly.
% The service handles printed text, hand-printed characters, and mixed-orientation
% layouts. The returned text is used verbatim as input to Stage~2; no additional
% image pre-processing is applied.

% \paragraph{Acoustic channel.}
% Audio inputs are transcribed using PhoWhisper-base~\cite{le2024phowhisper},
% served locally via the \texttt{faster-whisper} library with
% \texttt{compute\_type=int8}. The model operates on 16 kHz mono WAV files;
% stereo inputs are down-mixed before inference. Following transcription, a
% post-correction step is applied: a rule table of 47 domain-specific phonetic
% substitutions (e.g., Vietnamese homophone pairs for common electrical product
% names) is applied via exact string replacement. The substitution table was
% constructed by manual inspection of transcription errors on a development
% set of 120 audio samples.

% \subsection{Stage 2: LLM-Based Entity Extraction}
% \label{sec:extraction}

% \textbf{Input}: Transcribed text string $t$ from Stage~1.\\
% \textbf{Output}: A JSON object $\mathcal{E}$ conforming to a fixed schema
% with fields \texttt{merchant}, \texttt{line\_items} (array of
% \{\texttt{product}, \texttt{quantity}, \texttt{unit}\}), and
% \texttt{confidence}.

% The extraction prompt follows a four-component structure:
% \begin{enumerate}
%     \item \textbf{System instruction}: Defines the extraction task, target
%     entity schema, and output format constraints.
%     \item \textbf{Domain context}: A static block containing the phonetic
%     correction dictionary and ten common Vietnamese business abbreviations.
%     \item \textbf{Few-shot examples}: Three input–output pairs drawn from
%     distinct document types (handwritten invoice, audio transcript, printed
%     receipt).
%     \item \textbf{User turn}: The transcribed text $t$.
% \end{enumerate}

% The model (Gemini 1.5 Flash) is called with temperature $T = 0.2$ and JSON
% mode enforced. Temperature $T = 0.1$ is used for the taxonomy prediction
% sub-task (Section~\ref{sec:taxonomy}), where strict adherence to a fixed
% label set is required. JSON schema validation is applied to every response;
% malformed outputs trigger one retry before the record is routed to the HITL
% queue.

% Algorithm~\ref{alg:extraction} summarises the extraction procedure.

% \begin{algorithm}
% \caption{LLM-Based Entity Extraction}\label{alg:extraction}
% \begin{algorithmic}[1]
% \Require Transcribed text $t$, prompt template $P$, schema $S$, max retries $R=1$
% \Ensure Structured entity set $\mathcal{E}$ or HITL flag
% \State $prompt \gets \text{Build}(P, t)$
% \For{$attempt \gets 1$ to $R+1$}
%     \State $response \gets \text{LLM}(prompt,\; T=0.2,\; \text{json\_mode}=\texttt{true})$
%     \If{$\text{Validate}(response, S)$}
%         \State \Return $\text{Parse}(response)$
%     \EndIf
% \EndFor
% \State \Return $\textsc{HitlFlag}$ \Comment{Route to human review}
% \end{algorithmic}
% \end{algorithm}

% \subsection{Merchant Identity Resolution}
% \label{sec:merchant}

% \textbf{Input}: Raw merchant string $m$ extracted by $f_{\mathrm{IE}}$.\\
% \textbf{Output}: Canonical merchant identifier $d_m \in \mathcal{D}$, or
% exception flag.

% Merchant names in Vietnamese enterprise data are frequently subject to:
% (i) typographical errors from manual entry, (ii) abbreviation of legal
% entity suffixes (\textit{cty} $\to$ \textit{cong ty};
% \textit{tnhh} $\to$ \textit{trach nhiem huu han};
% \textit{ch} $\to$ \textit{cua hang}), and (iii) word-order variation in
% spoken references. Because merchant names are proper nouns, dense vector
% embeddings tend to conflate legally distinct entities that share generic
% vocabulary (e.g., \textit{cty tnhh thuong mai A} vs.\
% \textit{cty tnhh thuong mai B}).

% Resolution proceeds in two steps. First, normalisation applies lowercasing,
% whitespace collapsing, and abbreviation expansion via a 28-entry dictionary.
% Second, the normalised string is matched against the merchant registry using
% \texttt{token\_set\_ratio} from the RapidFuzz library, which computes:
% \begin{equation}
%     \text{TSR}(q, r) = \text{Levenshtein}\!\left(
%         \text{sort}(q \cap r),\;
%         \text{sort}(q \cap r) + \text{sort}(q \setminus r),\;
%         \text{sort}(q \cap r) + \text{sort}(r \setminus q)
%     \right) \times 100,
% \end{equation}
% where $q$ and $r$ are token multisets of the query and registry entry,
% respectively. A match is accepted when $\text{TSR}(q, r) \geq \tau_m = 80$;
% candidates below this threshold are routed to an exception queue.

% The threshold $\tau_m$ was selected by evaluating precision and recall on
% a development split of 200 labelled merchant strings at thresholds in
% $\{70, 75, 80, 85, 90\}$.

% \subsection{Product Search via Hybrid Retrieval}
% \label{sec:product}

% \textbf{Input}: Normalised product description string $p$ from $\mathcal{E}$.\\
% \textbf{Output}: Ranked list of SKU candidates with fusion scores.

% Product descriptions require both semantic understanding (to handle
% paraphrases and translations) and lexical precision (to distinguish, for
% example, \texttt{LED 8W} from \texttt{LED 9W}). We execute two retrieval
% streams in parallel:

% \paragraph{Dense retrieval.}
% Product descriptions are embedded using \texttt{keepitreal/vietnamese-sbert}
% (PhoBERT backbone~\cite{nguyen2020phobert}, 768-dimensional output,
% cosine similarity). The top-100 candidates are retrieved from the Qdrant
% collection by approximate nearest-neighbour search (HNSW index,
% $\texttt{ef} = 128$).

% \paragraph{Sparse retrieval.}
% Qdrant's built-in sparse vector index is populated with BM25-weighted
% term vectors~\cite{robertson2009probabilistic} ($k_1 = 1.5$, $b = 0.75$).
% The top-100 candidates are retrieved by inner product over sparse vectors.

% \paragraph{Fusion.}
% The two ranked lists are combined using \emph{Reciprocal Rank Fusion} (RRF)
% as the primary strategy~\cite{cormack2009reciprocal}:
% \begin{equation}
%     \text{RRF}(d) = \sum_{m \in \{\text{dense},\,\text{sparse}\}}
%         \frac{1}{k + r_m(d)}, \qquad k = 60.
%     \label{eq:rrf}
% \end{equation}
% Two additional fusion variants are evaluated in ablation:

% \emph{Weighted Score Fusion} (WSF) applies min-max normalisation to each
% stream's raw scores before a weighted sum:
% \begin{equation}
%     \text{WSF}(d) = \alpha \cdot \hat{s}_{\text{dense}}(d)
%                   + (1-\alpha) \cdot \hat{s}_{\text{sparse}}(d),
%     \label{eq:wsf}
% \end{equation}

% \emph{Weighted RRF} (WRRF) applies $\alpha$ directly to the rank-based
% contributions, avoiding score normalisation:
% \begin{equation}
%     \text{WRRF}(d) = \frac{\alpha}{k + r_{\text{dense}}(d)}
%                    + \frac{1-\alpha}{k + r_{\text{sparse}}(d)}.
%     \label{eq:wrrf}
% \end{equation}

% When $\alpha = 0.5$, WRRF reduces to standard RRF (Equation~\ref{eq:rrf}).
% All three strategies are compared in Section~\ref{sec:ablation}.

% \subsection{Taxonomy-Guided Retrieval Filter}
% \label{sec:taxonomy}

% For catalogs exceeding a configurable size threshold (default: 10{,}000
% active SKUs) or for categories with high inter-class embedding similarity,
% an optional pre-filtering step constrains the retrieval search space.

% \textbf{Input}: Product description $p$.\\
% \textbf{Output}: Taxonomy label tuple $(L_1, L_2, L_3)$ and a filtered
% Qdrant payload filter.

% The LLM (Gemini 1.5 Flash, $T = 0.1$) classifies $p$ into the three-level
% enterprise taxonomy. The predicted labels are applied as a \texttt{must}
% payload filter in Qdrant, reducing the candidate pool before both dense and
% sparse retrieval execute. This filter is disabled when the catalog has
% fewer than the threshold or when the LLM classification confidence (measured
% by consistency across three independent samples at $T=0$) falls below 0.7.

% \subsection{Human-in-the-Loop Fallback}

% A confidence threshold $\tau_p = 0.015$ is applied to the top-1 RRF score
% after fusion. Given $k = 60$, this threshold is equivalent to requiring the
% top candidate to be ranked within position~6 by at least one retrieval
% stream ($\frac{1}{60+6} \approx 0.015$). Candidates below this threshold,
% and all merchant resolution failures, are routed to a human review interface
% that presents the top-$K$ ($K = 5$) candidates alongside the original input
% for manual confirmation before database write.

% \subsection{Prompt Design}
% \label{sec:prompt}

% Table~\ref{tab:prompt} lists the prompt parameters used throughout
% the pipeline. All prompts are version-controlled and included in the
% supplementary material.

% \begin{table}[htbp]
% \centering
% \caption{Prompt and inference parameters for all LLM invocations.}
% \label{tab:prompt}
% \begin{tabular}{llll}
% \toprule
% \textbf{Task} & \textbf{Model} & \textbf{Temperature} & \textbf{Output format} \\
% \midrule
% Entity extraction   & Gemini 1.5 Flash & 0.2 & JSON (schema-enforced) \\
% Taxonomy prediction & Gemini 1.5 Flash & 0.1 & JSON (fixed label set)  \\
% ASR post-correction & Rule table       & n/a & Plain text              \\
% \bottomrule
% \end{tabular}
% \end{table}

% \subsection{Implementation and Reproducibility}
% \label{sec:reproducibility}

% \paragraph{Hardware.}
% All experiments are conducted on a single server with an Intel Xeon E5-2690
% CPU (16 cores), 64 GB RAM, and one NVIDIA A10G GPU (24 GB VRAM).
% PhoWhisper-base inference runs on the GPU; all other components run on CPU.

% \paragraph{Software.}
% Python 3.11; \texttt{faster-whisper} 1.0.3 (\texttt{int8} quantisation);
% Qdrant 1.9.2 (local Docker deployment); \texttt{rapidfuzz} 3.6.1;
% \texttt{sentence-transformers} 2.7.0; Google Generative AI Python SDK 0.5.4.

% \paragraph{Model identifiers.}
% \begin{itemize}
%     \item ASR: \texttt{vinai/PhoWhisper-base} (HuggingFace Hub)
%     \item Embeddings: \texttt{keepitreal/vietnamese-sbert} (HuggingFace Hub)
%     \item LLM: \texttt{gemini-1.5-flash-001} (Google AI API)
% \end{itemize}

% \paragraph{Hyperparameters.}
% Table~\ref{tab:hyperparams} lists all tunable parameters, their values, and
% the selection method.

% \begin{table}[htbp]
% \centering
% \caption{System hyperparameters, values, and selection method.}
% \label{tab:hyperparams}
% \begin{tabular}{llll}
% \toprule
% \textbf{Parameter} & \textbf{Symbol} & \textbf{Value} & \textbf{Selection} \\
% \midrule
% RRF smoothing constant        & $k$         & 60   & Standard IR practice~\cite{cormack2009reciprocal} \\
% Retrieval pool size (each)    & —           & 100  & Fixed \\
% Fusion weight                 & $\alpha$    & 0.5  & Grid search on dev set \\
% Merchant match threshold      & $\tau_m$    & 80\% & Precision–recall dev set \\
% HITL confidence threshold     & $\tau_p$    & 0.015 & Derived from $k=60$ \\
% Taxonomy filter threshold     & —           & 10{,}000 SKUs & Fixed \\
% LLM temperature (extraction)  & $T$         & 0.2  & Fixed \\
% LLM temperature (taxonomy)    & $T$         & 0.1  & Fixed \\
% \bottomrule
% \end{tabular}
% \end{table}

% \paragraph{Dataset splits.}
% The evaluation corpus (described in Section~\ref{sec:dataset}) is partitioned
% into development (15\%), test (15\%), and remaining training-equivalent
% documents used for threshold selection and phonetic substitution table
% construction. No test-set documents were seen during threshold selection.

% \paragraph{Prompt artefacts.}
% The full text of all prompts, few-shot examples, and the ASR phonetic
% substitution table are provided in the supplementary material to support
% replication.